\uuid{sjlA}
\exo7id{3499}
\titre{exo7 3499}
\auteur{quercia}
\organisation{exo7}
\datecreate{2010-03-10}
\isIndication{false}
\isCorrection{true}
\chapitre{Réduction d'endomorphisme, polynôme annulateur}
\sousChapitre{Valeur propre, vecteur propre}
\cpge{mpsi-mp2i.determinants.d, pcsi.calcul-matriciel-systemes-lineaires.d, ptsi.calcul-matriciel-systemes-lineaires.d, mpsi-mp2i.polynomes-fractions-rationnelles.c, pcsi.polynomes.c, ptsi.polynomes.c, bcpst1.matrices, bcpst1.polynomes-reels.a, mp-mpi.reduction-endomorphismes.b, pc.reduction.a, psi.reduction-endomorphismes.a, pt.reduction-endomorphismes.a, mp-mpi.reduction-endomorphismes.c, pc.reduction.b, psi.reduction-endomorphismes.b, pt.reduction-endomorphismes.b, bcpst2.valeurs-propres-vecteurs-propres.a, bcpst2.valeurs-propres-vecteurs-propres.b}
\module{Algèbre}
\niveau{L2}
\difficulte{}

\contenu{
\texte{
Chercher les valeurs propres des matrices :
}
\begin{enumerate}
    \item \question{$\begin{pmatrix}0      &\dots  &0      &1      \cr
          \vdots &       &\vdots &\vdots \cr
          0      &\dots  &0      &n-1    \cr
          1      &\dots  &n-1    &n      \cr \end{pmatrix}$.}
\reponse{$0$ et les racines de $6\lambda^2 -6n\lambda -n(n-1)(2n-1) = 0$.}
    \item \question{$\begin{pmatrix}0            &\sin\alpha &\sin2\alpha \cr
          \sin\alpha   &0          &\sin2\alpha \cr
          \sin2\alpha  &\sin\alpha &0           \end{pmatrix}$.}
\reponse{$\sin\alpha+\sin2\alpha,\quad -\sin\alpha,\quad -\sin2\alpha$.}
\end{enumerate}
}
